% =====================================================================
%  1bat-ud4-4.tex  —  SESSIÓ 4: Desviació típica i variància
%  (sense preàmbul: s'inclou des de main.tex)
% =====================================================================
\horatitol{Sessió 4 — Desviació típica i variància}%
{La xifra que captura la desigualtat: com de lluny estan, de mitjana, les dades respecte del centre.}

\subsection{Idea clau}

A la sessió 3 vam veure que dos edificis amb la \textbf{mateixa mitjana}
($\num{1500}\eur$) poden tenir realitats oposades. Va quedar pendent una pregunta:
\emph{com mesurem «com d'escampades» estan les dades?} Avui la responem amb la
\textbf{variància} i la \textbf{desviació típica}.

\begin{idea}
\textbf{Variància i desviació típica}\\
Mesuren \textbf{com de lluny} estan les dades respecte de la mitjana $\overline{x}$.
\begin{itemize}[leftmargin=1.4em, itemsep=3pt]
  \item \textbf{Variància} $\sigma^2$ — la mitjana de les \emph{distàncies al quadrat}
        a la mitjana:
        \[
        \sigma^2=\frac{(x_1-\overline{x})^2+\cdots+(x_N-\overline{x})^2}{N}
        =\frac{\sum (x_i-\overline{x})^2}{N}.
        \]
  \item \textbf{Desviació típica} $\sigma$ — la seva \emph{arrel quadrada}, que torna
        a les unitats originals:
        \[
        \sigma=\sqrt{\sigma^2}.
        \]
\end{itemize}
\textbf{Com es llegeix:} $\sigma$ \textbf{petita} $\rightarrow$ dades a prop de la
mitjana (grup \emph{homogeni}); $\sigma$ \textbf{gran} $\rightarrow$ dades molt
escampades (grup \emph{desigual}). \textbf{No t'oblidis l'arrel} al final!
\end{idea}

\subsection{Exemples}

\begin{exemple}[la desviació típica de l'edifici A]
A l'edifici A teníem (simplifiquem a cinc dades representatives per veure el
procediment): $\num{1400},\ \num{1450},\ \num{1500},\ \num{1550},\ \num{1600}$, amb
mitjana $\overline{x}=\num{1500}\eur$. Calculem les distàncies a la mitjana, les
elevem al quadrat i en fem la mitjana:
\[
\begin{aligned}
\sigma^2&=\frac{(-100)^2+(-50)^2+0^2+50^2+100^2}{5}
=\frac{\num{10000}+\num{2500}+0+\num{2500}+\num{10000}}{5}\\[2pt]
&=\frac{\num{25000}}{5}=\num{5000}.
\end{aligned}
\]
I la desviació típica és l'arrel: $\sigma=\sqrt{\num{5000}}\approx 70,7\eur$. És una
$\sigma$ \textbf{petita} comparada amb la mitjana: l'edifici A és \textbf{homogeni}.
\end{exemple}

\begin{exemple}[la desviació típica de l'edifici B]
A l'edifici B teníem vuit zeros i dos valors de $\num{7500}\eur$, amb la mateixa
mitjana $\overline{x}=\num{1500}\eur$. Les distàncies al quadrat:
\[
\sigma^2=\frac{8\cdot(0-\num{1500})^2+2\cdot(\num{7500}-\num{1500})^2}{10}
=\frac{8\cdot\num{2250000}+2\cdot\num{36000000}}{10}.
\]
\[
\sigma^2=\frac{\num{18000000}+\num{72000000}}{10}=\frac{\num{90000000}}{10}=\num{9000000}.
\]
\[
\sigma=\sqrt{\num{9000000}}\approx \num{3000}\eur.
\]
Una $\sigma$ \textbf{enorme} ($\approx\num{3000}\eur$, el doble de la mitjana!). Tot i
tenir la mateixa mitjana que A, la desviació típica \textbf{captura la desigualtat}
que la mitjana amagava: l'edifici B és \textbf{extremadament desigual}.
\end{exemple}

\subsection{Exercicis resolts}

\begin{exresolt}[desviació típica pas a pas]
Cinc voluntaris han fet aquestes hores: $4$, $6$, $8$, $6$, $6$. Calcula la mitjana,
la variància i la desviació típica.
\solucio
\textbf{Pas 1 — mitjana:}
\[
\overline{x}=\frac{4+6+8+6+6}{5}=\frac{30}{5}=6.
\]
\textbf{Pas 2 — distàncies al quadrat} (cada dada menys $6$, al quadrat):
\[
(4-6)^2=4,\quad (6-6)^2=0,\quad (8-6)^2=4,\quad (6-6)^2=0,\quad (6-6)^2=0.
\]
\textbf{Pas 3 — variància} (en fem la mitjana):
\[
\sigma^2=\frac{4+0+4+0+0}{5}=\frac{8}{5}=1,6.
\]
\textbf{Pas 4 — desviació típica} (arrel):
\[
\sigma=\sqrt{1,6}\approx 1,26\ \text{hores}.
\]
És petita: el grup és bastant homogeni (gairebé tots fan $6$ hores).
\resposta{$\overline{x}=6$; $\sigma^2=1,6$; $\sigma\approx 1,26$ hores.}
\end{exresolt}

\begin{exresolt}[comparar la dispersió de dos grups]
Dos grups d'usuaris esperen, en minuts, per ser atesos:
\[
\text{Grup 1: } 10,\ 12,\ 11,\ 9,\ 8 \qquad
\text{Grup 2: } 2,\ 20,\ 1,\ 18,\ 9.
\]
Comprova que tenen la mateixa mitjana i digues, calculant $\sigma$, quin grup té una
espera més \textbf{regular}.
\solucio
\textbf{Mitjanes:} $\overline{x}_1=\dfrac{10+12+11+9+8}{5}=\dfrac{50}{5}=10$ i
$\overline{x}_2=\dfrac{2+20+1+18+9}{5}=\dfrac{50}{5}=10$. Iguals.

\textbf{Grup 1} (distàncies al quadrat respecte de $10$): $0,4,1,1,4$, suma $10$:
\[
\sigma_1^2=\frac{10}{5}=2,\qquad \sigma_1=\sqrt{2}\approx 1,41\ \text{min}.
\]
\textbf{Grup 2} (distàncies a $10$): $(-8)^2,10^2,(-9)^2,8^2,(-1)^2=64,100,81,64,1$,
suma $310$:
\[
\sigma_2^2=\frac{310}{5}=62,\qquad \sigma_2=\sqrt{62}\approx 7,87\ \text{min}.
\]
El \textbf{grup 1} té una espera molt més regular ($\sigma\approx 1,41$) que el grup
2 ($\sigma\approx 7,87$), tot i la mateixa mitjana.
\resposta{Mateixa mitjana ($10$ min); el grup 1 és més regular ($\sigma_1\approx 1,41$
contra $\sigma_2\approx 7,87$ min).}
\end{exresolt}

\subsection{Exercicis per fer}

\begin{exercicis}
\begin{exlist}
  \item Calcula la \textbf{mitjana} i, després, la \textbf{variància} i la
        \textbf{desviació típica} d'aquest conjunt: $5$, $7$, $9$, $7$, $2$.
  \item Per al conjunt $20$, $20$, $20$, $20$ (totes les dades iguals), calcula
        $\sigma$. Quant val? Per què té sentit aquest resultat?
  \item Dos casals tenen aquestes inscripcions per setmana:
        \[
        \text{Casal P: } 30,\ 32,\ 31,\ 29,\ 28 \qquad
        \text{Casal Q: } 10,\ 50,\ 30,\ 5,\ 55.
        \]
        Comprova que tenen la mateixa mitjana i calcula la $\sigma$ de cadascun. Quin
        té les inscripcions més estables?
  \item \textbf{(Compte amb el procés)} Un company calcula la desviació típica d'un
        conjunt i dóna com a resposta $\sigma^2=36$. Quin pas li ha faltat? Quina és
        la desviació típica correcta?
  \item \textbf{(Interpretació)} Dos barris tenen la mateixa renda mitjana
        ($\num{20000}\eur$), però el barri A té $\sigma=\num{2000}\eur$ i el barri B
        té $\sigma=\num{12000}\eur$. Quin barri és més desigual? En quin creus que la
        intervenció social és més urgent i per què?
  \item \textbf{(Raonament)} Explica, amb les teves paraules, per què la desviació
        típica «captura» la desigualtat que la mitjana amaga. Per què elevem les
        distàncies \textbf{al quadrat} en comptes de sumar-les tal qual?
\end{exlist}
\end{exercicis}

% ---------------------------------------------------------------------
\fullsolucions{Sessió 4}
\begin{solucions}
\begin{sollist}
  \item $\overline{x}=\dfrac{5+7+9+7+2}{5}=\dfrac{30}{5}=6$. Distàncies al quadrat:
        $1,1,9,1,16$, suma $28$. $\sigma^2=\dfrac{28}{5}=5,6$;
        $\sigma=\sqrt{5,6}\approx 2,37$.
  \item Totes les dades són $20$, així que la mitjana és $20$ i totes les distàncies
        valen $0$: $\sigma^2=0$ i $\sigma=0$. Té sentit: si \textbf{no hi ha cap
        dispersió} (totes iguals), la desviació típica és nul·la.
  \item Mitjanes: $\overline{x}_P=\dfrac{150}{5}=30$ i $\overline{x}_Q=\dfrac{150}{5}=30$,
        iguals. Casal P: distàncies al quadrat $0,4,1,1,4$, suma $10$;
        $\sigma_P^2=2$, $\sigma_P\approx 1,41$. Casal Q: distàncies a $30$:
        $400,400,0,625,625$, suma $\num{2050}$; $\sigma_Q^2=410$,
        $\sigma_Q\approx 20,25$. El \textbf{casal P} té les inscripcions molt més
        estables.
  \item Li ha faltat l'\textbf{últim pas: l'arrel quadrada}. El que ha calculat ($36$)
        és la variància $\sigma^2$, no la desviació típica. La correcta és
        $\sigma=\sqrt{36}=6$.
  \item El \textbf{barri B} és molt més desigual ($\sigma=\num{12000}\eur$ contra
        $\num{2000}\eur$): tot i la mateixa renda mitjana, al barri B hi ha rendes
        molt altes i molt baixes barrejades. La intervenció és més urgent al barri B,
        perquè la gran dispersió indica que hi conviuen famílies en situacions molt
        diferents (algunes probablement molt vulnerables).
  \item Resposta oberta. Idees: la $\sigma$ mesura com de lluny estan les dades del
        centre, així que distingeix un grup homogeni d'un de desigual encara que
        tinguin la mateixa mitjana. Elevem al quadrat perquè, si sumàssim les
        distàncies amb el seu signe, les positives i les negatives es
        \textbf{cancel·larien} i sempre donaria zero; el quadrat les fa totes
        positives (i després l'arrel torna a les unitats originals).
\end{sollist}
\end{solucions}
